\section{Perbandingan Top-Down dan Bottom-Up}
\label{sec:ch11-perbandingan}

\subsection{Perbandingan}

\begin{table}[H]
\centering
\caption{Perbandingan parsing top-down dan bottom-up}
\label{tab:ch11-perbandingan}
\begin{tabularx}{\textwidth}{@{}lXX@{}}
\toprule
\textbf{Aspek} & \textbf{Top-down} & \textbf{Bottom-up} \\
\midrule
Arah & Akar ke daun & Daun ke akar \\
Left recursion & Harus dihapus & Dapat ditangani \\
Grammar & LL(1) dan variannya & LR, SLR, LALR, precedence \\
Tabel & Satu entri per nonterminal & Satu entri per state \\
Pesan kesalahan & Mudah dibuat kontekstual & Lebih sulit \\
Implementasi & Recursive descent & Stack dan tabel \\
Contoh & Parser LL(1) & Bison, Yacc, operator precedence \\
\bottomrule
\end{tabularx}
\end{table}

\subsection{Kapan Memilih Top-Down}

Top-down cocok ketika:
\begin{itemize}[leftmargin=*]
  \item grammar relatif sederhana dan dapat difaktorkan;
  \item pesan kesalahan yang baik menjadi prioritas;
  \item parser ditulis tangan untuk kontrol penuh;
  \item pembelajaran konsep parsing menjadi tujuan.
\end{itemize}

\subsection{Kapan Memilih Bottom-Up}

Bottom-up cocok ketika:
\begin{itemize}[leftmargin=*]
  \item grammar besar dan kompleks;
  \item left recursion alami dan tidak ingin diubah;
  \item parser generator tersedia;
  \item kinerja dan cakupan grammar lebih penting daripada
        kesederhanaan implementasi.
\end{itemize}

\subsection{Jembatan ke LR}

Operator precedence adalah bentuk sederhana dari parsing bottom-up.
Bab XII akan memperkenalkan LR parsing yang lebih umum: item LR(0),
SLR(1), canonical LR(1), dan LALR(1). LR menggabungkan kekuatan
bottom-up dengan tabel yang dapat dibangun secara sistematis.

\begin{infobox}
Banyak compiler produksi memakai parser generator LR karena mampu
menangani grammar bahasa nyata tanpa banyak transformasi. Namun,
memahami shift-reduce dan precedence tetap penting untuk membaca
konflik dan pesan diagnostik yang dihasilkan generator.
\end{infobox}
